\section{Auxiliary proofs}

\begin{lemma}[Logarithmic linearization]
For \(a>0\) and \(x>0\),
\[
 \log x\le ax+\log(1/a)-1.
\]
\end{lemma}
\begin{proof}
The function \(\log x-ax\) reaches its global maximum at \(x=1/a\), where
its value is \(\log(1/a)-1\).
\end{proof}

\begin{lemma}[Shifted Gaussian lattice]
For \(\alpha>0\) and \(c\in\mathbb R\),
\[
 \sum_{q\in\mathbb Z}\ee^{-\alpha(q-c)^2}
 \le 1+\sqrt{\frac\pi\alpha}.
\]
\end{lemma}
\begin{proof}
Poisson summation shows that the shifted sum is maximized when \(c\) is an
integer, because every Fourier coefficient of the Gaussian is positive.
At an integral shift,
\[
 \sum_{q\in\mathbb Z}\ee^{-\alpha q^2}
 =1+2\sum_{q\ge1}\ee^{-\alpha q^2}.
\]
Since \(x\mapsto\ee^{-\alpha x^2}\) decreases on \([0,\infty)\), each term
in the positive tail is bounded by the integral over the preceding unit
interval.  Thus
\[
 2\sum_{q\ge1}\ee^{-\alpha q^2}
 \le2\int_0^\infty\ee^{-\alpha x^2}\,dx
 =\sqrt{\frac\pi\alpha}.
\]
\end{proof}

\begin{proposition}[Coefficient summation]
Suppose \(a_0=1\) and for \(q\ge1\),
\[
 |a_q|\le\exp(-\alpha q^2+Bq),
 \qquad \alpha>0.
\]
Then
\[
 \left|\sum_{q\ge0}a_q\right|
 \le\left(1+\sqrt{\frac\pi\alpha}\right)
 \exp\left(\frac{B^2}{4\alpha}\right).
\]
\end{proposition}
\begin{proof}
Extend the positive majorant from \(q\ge0\) to all integer \(q\), complete
the square, and apply the preceding lemma with \(c=B/(2\alpha)\).
\end{proof}

The passage from finite Taylor truncations to \(\Dpol\) is inherited from the
order-zero nuclear theorem.  The present paper changes only the numerical
majorant and does not introduce a second determinant limit.

